\section{Orientation averages}\label{sec:averages}

\subsection{The rule of weighted means}\label{sec:averages-rule}

In an untextured polycrystal the grain orientations are uniformly distributed, and the
average of an element of a crystal algebra over the rotations,
$\avg{\tK}=\int_{\SOrth}(Q\star)\,\tK\,(Q\star)^{\mathsf T}\,\dd Q$, is isotropic. The
average is linear, idempotent, self-adjoint for $\langle\cdot,\cdot\rangle$ and fixes
$\Alg{\clsiso}$, so it is the orthogonal projection onto $\Alg{\clsiso}$, and the last
formula of~\eqref{eq:block-rules} gives its moduli. By~\eqref{eq:cub-iso-units}
and~\eqref{eq:hex-iso-units}, each isotropic unit is a sum of crystal units; hence the
isotropic moduli of $\avg{\tK}$ are means of the crystal block entries weighted by the
dimensions of the irreducible components,
\begin{equation}
m_\ell=\sum_p\frac{\dimirr{p}}{2\ell+1}\,K^p_{c_\ell c_\ell}\quad(\ell=3,2),\qquad
M_{ab}=\sum_p\frac{\dimirr{p}}{3}\,K^p_{ab}\quad(a,b\in\{s,r\}),\label{eq:weighted-means}
\end{equation}
where $c_\ell$ is the copy of $\Hrm{\ell}$ in $p$ and the sums run over the $p$ that
contain one. The weights add up to one, the couplings between copies of different
$\Hrm{\ell}$ do not survive, and the rule holds for the non-symmetric elements too. For
the Voigt estimate $\Vg{\tD}=\avg{\tD}$ of a cubic tensor,
\begin{equation}
\Vg{m_3}=\tfrac17\big(3M_{hh}+3N_{33}+m_A\big),\qquad
\Vg{m_2}=\tfrac15\big(3N_{22}+2m_E\big),\qquad \Vg{M}=M_{(sr)},\label{eq:cub-voigt}
\end{equation}
with $M_{(sr)}$ the $(s,r)$ part of $M$. The Reuss estimate
$\Rs{\tD}=\avg{\tD^{-1}}^{-1}$ applies the rule to the inverse blocks and inverts each
isotropic block; with $\bm m_h=(M_{hs},M_{hr})^{\mathsf T}$,
\begin{equation}
\begin{gathered}
\frac1{\Rs{m_3}}=\frac17\Big(3(M^{-1})_{hh}+3(N^{-1})_{33}+\frac1{m_A}\Big),\qquad
\frac1{\Rs{m_2}}=\frac15\Big(3(N^{-1})_{22}+\frac2{m_E}\Big),\\
\Rs{M}=M_{(sr)}-\frac{\bm m_h\bm m_h^{\mathsf T}}{M_{hh}},
\end{gathered}\label{eq:cub-reuss}
\end{equation}
the last being the Schur complement of $M_{hh}$, the inverse of $(M^{-1})_{(sr)}$. For a
hexagonal tensor,
\begin{equation}
\begin{gathered}
\Vg{m_3}=\tfrac17\big(\hexA_{33}+m_{B_1}+m_{B_2}+2\hexF_{33}+2\hexG_{33}\big),\qquad
\Vg{m_2}=\tfrac15\big(m_{A_1}+2\hexF_{22}+2\hexG_{22}\big),\\
\Vg{M}=\tfrac13\big(\hexA_{(sr)}+2\hexF_{(sr)}\big),\qquad
(\Rs{M})^{-1}=\tfrac13\big((\hexA^{-1})_{(sr)}+2(\hexF^{-1})_{(sr)}\big),\\
\frac1{\Rs{m_3}}=\frac17\Big(c+\frac1{m_{B_1}}+\frac1{m_{B_2}}\Big),\qquad
\frac1{\Rs{m_2}}=\frac15\Big(\frac1{m_{A_1}}+2(\hexF^{-1})_{22}+2(\hexG^{-1})_{22}\Big),
\end{gathered}\label{eq:hex-voigt-reuss}
\end{equation}
with $c=(\hexA^{-1})_{33}+2(\hexF^{-1})_{33}+2(\hexG^{-1})_{33}$. The Hill estimate is
$\Hl{\tD}=\tfrac12(\Vg{\tD}+\Rs{\tD})$. In the constants $a_i$, \eqref{eq:cub-voigt}
reproduces the Voigt average of Lazar et al.~\cite[eq.~(158)]{Lazar2022}. Classical closed
forms are obtained from rotation invariants, for third-order elastic
constants~\cite[eqs.~(2.4) and~(2.6)]{Blaschke2017} and for photoelastic tensors, which
lack the major symmetry~\cite{Ranganath1973}. Kube and Norris~\cite{Kube2026} obtain the
Voigt and Reuss averages of cubic third-order constants in a basis of Walpole's type; here
the blocks give both for the gradient tensor, without invariants.

\subsection{Reuss below Voigt}\label{sec:averages-order}

The Reuss estimate needs only $\tD$ and $\avg{\tD^{-1}}$ to be invertible, but it is the
response of the grains under a uniform double stress only if $\tD\succ0$. We use it, and
the Hill estimate, only then, which excludes the Fe and W data
(section~\ref{sec:compliance-data}).

\emph{Proposition.} If $\tD\succ0$, then $\Rs{\tD}\preceq\Vg{\tD}$, that is,
$\Rs{m_3}\le\Vg{m_3}$, $\Rs{m_2}\le\Vg{m_2}$ and $\Rs{M}\preceq\Vg{M}$, with
$\Rs{\tD}=\Vg{\tD}$ if and only if $\tD$ is isotropic.

The orientation average is a positive unital linear map, so the inequality is Choi's,
$\avg{\tD}^{-1}\preceq\avg{\tD^{-1}}$~\cite[theorem~2.3.6]{Bhatia2007}; we give the short
proof because it also yields the case of equality.

\emph{Proof.} Let $\tD_Q=(Q\star)\tD(Q\star)^{\mathsf T}$. Since $\avg{\tD_Q}=\Vg{\tD}$ and
$\Vg{\tD}$ commutes with every $Q\star$, with $\mathbb V=(\Vg{\tD})^{-1}$,
\begin{equation}
\begin{aligned}
(\Rs{\tD})^{-1}-(\Vg{\tD})^{-1}
&=\avg{\tD_Q^{-1}-2\mathbb V+\mathbb V\tD_Q\mathbb V}\\
&=\avg{(\mathbb V-\tD_Q^{-1})\,\tD_Q\,(\mathbb V-\tD_Q^{-1})}\succeq0,
\end{aligned}
\label{eq:reuss-below-voigt}
\end{equation}
which is the Schur complement of the average of the positive semi-definite matrix
$\left(\begin{smallmatrix}\tD_Q&\tI\\\tI&\tD_Q^{-1}\end{smallmatrix}\right)$~\cite[theorem~1.3.3]{Bhatia2007}.
Inversion reverses the order of positive definite maps, so $\Rs{\tD}\preceq\Vg{\tD}$, and
since both are isotropic the order holds block by block. The integrand has the same trace for every
$Q$, positive unless $\tD=\Vg{\tD}$; hence equality holds only for isotropic $\tD$.
$\square$

In classical elasticity the Voigt and Reuss estimates bound the effective
stiffness~\cite{Blaschke2017,Kube2016}. Here they do not: they are the responses of the
grains under a uniform strain gradient or a uniform double stress. The heterogeneity
of $\tC$ alone induces gradient moduli that grow with the square of the grain
size~\cite{Malyi2021}, for which Voigt and Reuss bounds of the gradient effective medium
have been derived on spherical volume elements~\cite{Muhlich2012}; they are absent
from any average of $\tD$. For third-order elastic
constants, too, the bounds are lost~\cite{Blaschke2017}. We use them as algebraic
estimates of the part of the effective tensor carried by $\tD$, and the Hill mean as a
convention, empirical already in the classical case~\cite{Kube2016}.

\subsection{Cubic grains}\label{sec:averages-cubic}

By~\eqref{eq:cub-voigt} and~\eqref{eq:cub-reuss},
$\Vg{M}-\Rs{M}=\bm m_h\bm m_h^{\mathsf T}/M_{hh}$ has rank at most one: one generalised
eigenvalue of the pair $(\Rs{M},\Vg{M})$ equals one, and the two estimates agree on
$\Hrm1$ if and only if $M_{hs}=M_{hr}=0$. This continues a classical sequence. In a cubic
crystal the bulk part lies in a single irreducible block, so $\kappa_V=\kappa_R$, and
$\mu_V-\mu_R$ is proportional to the square of the anisotropy $c_{11}-c_{12}-2c_{44}$~\cite[eqs.~(27)
and~(28)]{Kube2016}. Here $\Hrm1$ lies in $\irr{T_1}$, of multiplicity three, and the
difference is of rank one. In the hexagonal class $\Hrm1$ enters $\irr{A_2}$ and $\irr{E_1}$
together with copies of $\Hrm3$ and $\Hrm2$, and the rank is two in general. A second
property follows from~\eqref{eq:cub-reuss}:
\begin{equation}
\frac{\Rs{m_2}}{m_E}=\frac{5/2}{1+\tfrac32\,m_E\,(N^{-1})_{22}}\;\to\;\frac52
\quad(m_E\to0),\label{eq:reuss-factor}
\end{equation}
where $\tfrac52=5/\dimirr{E}$: the nearly vanishing $\irr{E}$ part, two of the five
dimensions of $\Hrm2$, dominates the harmonic mean, as $\mu_R\to\tfrac52c'$ when
$c'=\tfrac12(c_{11}-c_{12})\to0$ in the classical case, by the Reuss modulus
of~\cite[eq.~(28)]{Kube2016} (their eq.~(30) prints the prefactor $\tfrac53$ for $5$).

Table~\ref{tab:averages} gives the estimates of the five data sets, and
figure~\ref{fig:cubic-averages}(a) the ratios of Reuss to Voigt. The Voigt estimates of Al
and W reproduce table~6 of Lazar et al.~\cite{Lazar2022} to $\NumTableSixMaxDiff$~eV/\AA.
For $\Hrm3$, the ratio is $\NumAlRVmthree$ (Al), $\NumCuMEAMRVmthree$ (Cu MEAM) and
$\NumCuEAMRVmthree$ (Cu EAM); for $\Hrm2$, $\NumAlRVmtwo$, $\NumCuMEAMRVmtwo^{\dagger}$ and
$\NumCuEAMRVmtwo$, the Cu MEAM value set by its small $m_E$, with
$\Rs{m_2}/m_E=\NumCuMEAMReussRatio$. The two copper tensors differ more than the Voigt and
Reuss estimates of either: the difference belongs to the potentials, not to copper.

These estimates must be weighed against the moduli induced by the heterogeneity of
$\tC$. Malyi~\cite[eqs.~(10.21) and (10.22)]{Malyi2021} gives them, to second order in the
anisotropy $c_{11}-c_{12}-2c_{44}$, for spherical grains of radius $R$ with independent
orientations, through his $L^2=2R^2/5$. Taken to form~II, they have the norm of the Voigt
estimate of $\tD$ at a grain diameter $2R$ of $\NumAlInducedCrossoverNm$~nm (Al),
$\NumCuMEAMInducedCrossoverNm$~nm (Cu MEAM) and $\NumCuEAMInducedCrossoverNm$~nm (Cu
EAM), and at $2R=100$~nm they exceed it by factors of about
$\NumAlInducedRatioHundredNm$ (Al) and $\NumCuMEAMInducedRatioHundredNm$ (Cu). This is an order of
magnitude, but it settles the scale: above a few nanometres the gradient response of a
polycrystal is set by the heterogeneity of $\tC$, and the averages of $\tD$ describe
the remainder.

For a physical reading we define block lengths $\lblk{p}$ by
\begin{equation}
\tD^p\bm w=\lblk{p}^2\,(\tC\ot\kron)^p\,\bm w,\label{eq:block-lengths}
\end{equation}
the generalised eigenvalues of each block of $\tD$ relative to that of the separable tensor
$\tC\ot\kron$; $\tD$ has the Helmholtz form $\ell^2\tC\ot\kron$ if and only if all
$\lblk{p}$ equal $\ell$. Each gradient estimate is paired with the estimate of $\tC$ of the
same kind. The pairing is a normalisation, not a joint estimate: a uniform strain and a
uniform strain gradient cannot be imposed together~\cite[remark~1]{Lazar2022}. Since $\avg{\tC\ot\kron}=\avg{\tC}\ot\kron$ and
$(\tC\ot\kron)^{-1}=\tC^{-1}\ot\kron$, a Helmholtz grain then has
$\lblk{p}=\ell$ in all three. The length of Lazar and Po~\cite[eq.~(107)]{Lazar2024} is
their dimension-weighted root mean square:
\begin{equation}
\lLP^2=\tfrac1{18}\,\big(\tC^{-1}\big)_{ijkl}D_{ijmklm}
=\tfrac1{18}\tr\big((\tC\ot\kron)^{-1}\tD\big)=\tfrac1{18}\sum_p\dimirr{p}\sum\lblk{p}^2,
\label{eq:lazar-po-length}
\end{equation}
the inner sum running over the lengths of block $p$. In the Voigt estimates
(table~\ref{tab:block-lengths}, figure~\ref{fig:cubic-averages}(b)) the block lengths
spread by factors $\NumAlVoigtLengthSpread$, $\NumCuMEAMVoigtLengthSpread$ and
$\NumCuEAMVoigtLengthSpread$, far from a single length; $\lLP$ of the Hill estimates is
$\NumAlHillLengthLP$, $\NumCuMEAMHillLengthLP$ and $\NumCuEAMHillLengthLP$~\AA. The Cu EAM
lengths of Reuss exceed those of Voigt because $\mu_R<\mu_V$ in the normalisation, without
contradicting the proposition.

\subsection{Hexagonal grains}\label{sec:averages-hexagonal}

By~\eqref{eq:hex-voigt-reuss} the Voigt estimate depends on $m_{B_1}+m_{B_2}$ only: by
\eqref{eq:hex-extra}, the extra modulus multiplies a harmonic tensor of degree six, whose
average vanishes. In the Reuss estimate only $\Rs{m_3}$ depends on it, through
$1/m_{B_1}+1/m_{B_2}=8\sigma/(4\sigma^2-\Dhex^2)$, $\sigma=\tfrac12(m_{B_1}+m_{B_2})$:
\begin{equation}
\Rs{m_3}(\Dhex)=\frac{7}{c+8\sigma/(4\sigma^2-\Dhex^2)},\label{eq:hex-reuss-delta}
\end{equation}
even in $\Dhex$, decreasing in $|\Dhex|$ and vanishing as $|\Dhex|\to2\sigma$, where
$\tP{B_1}$ or $\tP{B_2}$ loses positivity. No average depends on the sign of $\Dhex$,
that is, on the family of twofold axes chosen for $\bm a$.

No hexagonal gradient tensor of a crystal has been published, but granular
micromechanics gives one. For equal spheres with normal and tangential contact
stiffnesses $k_n$ and $k_s$, Pirmoradi et al.~\cite[eq.~(70)]{Pirmoradi2024} write the
sixth-order tensor, in form~I with the $\kappa$ of section~\ref{sec:form-II}, as an
integral over the contact directions weighted by their density, expanded in spherical
harmonics~\cite[eq.~(71)]{Pirmoradi2024}. They work out isotropic, transversely isotropic,
orthotropic, monoclinic and triclinic fabrics; the hexagonal one is our choice within their
eq.~(71). With the axial terms $a_{20}$, $a_{40}$, $a_{60}$ and the sixfold term
$a_{66}P_6^6(\cos\theta)\cos6\varphi$, where $P_6^6(x)=10395\,(1-x^2)^3$ is the
unnormalised associated Legendre function, the density is invariant under
the hexagonal group and even in the direction, so the fifth-order coupling vanishes, and
the classical tensor, which sees only the terms of degree up to four, is transversely
isotropic. In the basis, $a_{66}$ changes only $m_{B_1}$ and $m_{B_2}$:
\begin{equation}
\Dhex=\tfrac{1440}{13}\,(k_n-k_s)\,a_{66}\,\rho_cl^4,\label{eq:hex-granular}
\end{equation}
with $\rho_c$ the number of contacts per unit volume and $l$ the branch length; for
uniform contacts the tensor reproduces their eq.~(74). The extra modulus is thus the
sixfold harmonic of the contact distribution, weighted by the difference of the two
contact stiffnesses. A positive density bounds it: for $k_s=k_n/2$, $a_{20}=0.2$ and $a_{40}=a_{60}=0$, within
the range of their figure~3, $|\Dhex|\le\NumHexGranDeltaMax\,\sigma$, and $\Rs{m_3}$ changes
by at most $\NumHexGranReussChangePercent\,\%$. Sweeping $\Dhex$ further on the same base
(figure~\ref{fig:hexagonal}), $\Rs{m_3}(\sigma)/\Rs{m_3}(0)$ is $\NumHexTIReussAtSigma$, and
$21/23=\NumHexIsoReussAtSigma$ for every isotropic base, for which $c=5/m_3$ and
$\sigma=m_3$: the effect of the extra modulus is of second order.

\begin{table}[pos=tbp]
\caption{Voigt, Reuss and Hill estimates~\eqref{eq:cub-voigt}--\eqref{eq:cub-reuss} of the
cubic data sets of table~\ref{tab:block-moduli}: isotropic block moduli~\eqref{eq:iso-tensor}
(eV/\AA). Fe and W, which are not positive definite, have only the Voigt estimate. The last
digit shown is uncertain by at most five units; $\dagger$ marks a relative uncertainty
above 1\,\%, from the $m_E$ of Cu MEAM.}\label{tab:averages}
\centering\small
\input{averages}
\end{table}

\begin{table}[pos=tbp]
\caption{Block lengths~\eqref{eq:block-lengths} of the estimates (\AA), each relative to
$\tC\ot\kron$ with the classical stiffness of the same estimate, unlike
table~\ref{tab:application}, which uses the stiffness of Kr\"oner's scheme throughout;
$\lblk{1}^{\mp}$ are the two lengths of the $\Hrm1$ block, and
$\lLP$~\eqref{eq:lazar-po-length}. The row `crystal' gives $\lLP$ of the single crystal,
which for Al and Cu MEAM reproduces table~2 of Lazar and Po~\cite{Lazar2024}. Digits and
$\dagger$ as in table~\ref{tab:averages}.}\label{tab:block-lengths}
\centering\small
\input{block_lengths}
\end{table}

\begin{figure}[pos=tbp]
\centering
\includegraphics{cubic_averages}
\caption{Orientation averages of the positive definite cubic data. (a) Reuss over Voigt for
$m_3$ ($\Hrm3$), $m_2$ ($\Hrm2$) and, in the rows $\Hrm1^{\mp}$, the lower and upper
generalised eigenvalues of $(\Rs{M},\Vg{M})$; the upper one is exactly one. (b) Block
lengths of the Voigt (filled) and Reuss (open) estimates, as in
table~\ref{tab:block-lengths}, with $\lblk{1}^{\mp}$ in the rows $\Hrm1^{\mp}$, and $\lLP$ of
Voigt (dotted lines).}\label{fig:cubic-averages}
\end{figure}

\begin{figure}[pos=tbp]
\centering
\includegraphics{hexagonal_extra_modulus}
\caption{Reuss and Hill moduli $m_3$ over the Voigt modulus, against $\Dhex/\sigma$ with
$m_{B_1,B_2}=\sigma\pm\Dhex/2$, on two bases: any isotropic tensor, and the transversely
isotropic (TI) tensor of a granular packing with $k_s=k_n/2$ and
$a_{20}=0.2$~\cite{Pirmoradi2024}. Shaded: the values of $\Dhex$ that a positive contact
density gives to this packing, by~\eqref{eq:hex-granular}. Voigt does not depend on
$\Dhex$.}\label{fig:hexagonal}
\end{figure}
